\subsection*{Teil A: Zufallsexperiment oder nicht? (20 Minuten)}

Entscheide bei jeder Situation, ob ein Zufallsexperiment vorliegt. Kreuze an und begründe mit den vier Bedingungen aus dem Merkkasten.

\renewcommand{\arraystretch}{1.6}
\begin{center}
\begin{tabularx}{\textwidth}{|c|X|c|p{4.4cm}|}
\hline
\textbf{Nr.} & \textbf{Situation} & \textbf{Zufallsexp.?} & \textbf{Begründung} \\
\hline
A1 & Ein normaler Spielwürfel wird einmal geworfen. & $\square$ ja \; $\square$ nein & \\
\hline
A2 & Du rechnest $3 + 4$ aus. & $\square$ ja \; $\square$ nein & \\
\hline
A3 & Aus einem Beutel mit 3 roten und 2 blauen Kugeln wird ohne Hinsehen eine Kugel gezogen. & $\square$ ja \; $\square$ nein & \\
\hline
A4 & Ein Stein wird aus 2\,m Höhe losgelassen. Man beobachtet, ob er nach unten fällt. & $\square$ ja \; $\square$ nein & \\
\hline
A5 & Ein Glücksrad mit 4 gleich großen Feldern wird einmal gedreht. & $\square$ ja \; $\square$ nein & \\
\hline
A6 & Eine Münze wird mit der Zahl nach oben auf den Tisch gelegt. & $\square$ ja \; $\square$ nein & \\
\hline
A7 & Aus einem gut gemischten Kartenstapel wird eine Karte gezogen und ihre Farbe notiert. & $\square$ ja \; $\square$ nein & \\
\hline
A8 & Wasser wird auf 100\,°C erhitzt. Man beobachtet, ob es kocht. & $\square$ ja \; $\square$ nein & \\
\hline
A9 & Auf dem Schulfest zieht Mia ein Los aus der Lostrommel (Gewinn oder Niete). & $\square$ ja \; $\square$ nein & \\
\hline
A10 & Ben greift blind in eine Tüte mit roten, grünen und gelben Gummibärchen. & $\square$ ja \; $\square$ nein & \\
\hline
A11 & Lisa sucht sich morgens bewusst ihr Lieblings-T-Shirt aus dem Schrank aus. & $\square$ ja \; $\square$ nein & \\
\hline
A12 & Man beobachtet, ob es morgen an deinem Wohnort regnet. & $\square$ ja \; $\square$ nein & \\
\hline
\end{tabularx}
\end{center}
\renewcommand{\arraystretch}{1}
